% !TEX root = ../main.tex
\documentclass[../main.tex]{subfiles}
\begin{document}

\chapter{Software de visualización PC}
\label{chap:software}

La aplicación \texttt{vicon\_view} es un programa de escritorio desarrollado en C++ con Qt~6 (interfaz gráfica) y OpenCV, este último para la detección de rostros. Se comunica con la FPGA a través del chip FT232H mediante la API oficial de FTDI (\texttt{FTD2XX}), recibiendo el flujo de imagen en escala de grises y permitiendo el envío de comandos de control. El proyecto se compila con CMake (\texttt{CMakeLists.txt}),para la generación automática de los ficheros de Qt a partir de \texttt{mainwindow.ui} y \texttt{resources.qrc}.

La elección de C++ para el software de PC no fue casual o basada únicamente en el requisito, sino consecuencia directa de un problema detectado en una primera versión desarrollada en Python. Con Python, la recepción del vídeo fallaba de forma intermitente: la \ac{FIFO} asíncrona de la \ac{FPGA} se desbordaba (\texttt{overflow}) y se perdía el sincronismo de trama.

El diagnóstico con ILA reveló que la causa no era el ancho de banda —el enlace FT232H ofrece hasta 60~MB/s— sino la \textbf{latencia del host} en la lectura del buffer. El flujo de imagen es un caudal continuo: mientras el PC no vacía el buffer del FT232H, la \ac{FPGA} sigue escribiendo en la \ac{FIFO} asíncrona, que tiene una capacidad limitada. Si el PC deja de leer durante un intervalo demasiado largo, la \ac{FIFO} se llena y se produce el desbordamiento.

En Python, estos intervalos resultaban excesivos e impredecibles: se introducían pausas de varios milisegundos en la ejecución, durante las cuales no se leía del FT232H. Se midieron latencias de lectura del orden de $7$~ms, cuando la capacidad de la \ac{FIFO} solo permitía tolerar pausas de aproximadamente $2{,}6$~ms antes de desbordar. Es decir, las pausas superaban en más del doble el margen disponible.

La migración a C++ resolvió el problema. Al ejecutar la recepción en un hilo dedicado con alta prioridad, las latencias de lectura se reducen notablemente, dentro del margen que permite la \ac{FIFO}. Esto permite mantener la máxima tasa de fotogramas sin necesidad de recurrir a soluciones como reducir los \ac{FPS} o sobredimensionar la \ac{FIFO}, y confirma que el cuello de botella residía por completo en el lado del PC, no en la \ac{FPGA}.

\section{Arquitectura y concurrencia}
\label{sec:arquitectura_sw}

A diferencia de versiones anteriores basadas en una única llamada bloqueante a \texttt{FT\_Read} en el hilo principal (lo que llegaba a congelar la interfaz), la aplicación actual separa el trabajo en \textbf{tres hilos} independientes, comunicados con la \ac{GUI} exclusivamente mediante señales y \textit{slots} de Qt (nunca actualizando directamente un widget desde un hilo secundario):

\begin{itemize}
    \item \textbf{Hilo principal} (\texttt{QApplication::exec()}): ejecuta el bucle de eventos de Qt, gestiona los \textit{widgets} generados desde \texttt{mainwindow.ui} y actualiza la imagen, el indicador de FPS y el estado de conexión en los \textit{slots} \texttt{onFrameReady}, \texttt{onConnectionStatusChanged} y \texttt{onFacesDetected}.
    \item \textbf{Hilo de captura} (\texttt{captureThreadFunc}): en Windows se eleva a prioridad \texttt{THREAD\_PRIORITY\_TIME\_CRITICAL}. Llama a \texttt{FT\_Read} de forma directa (sin mutex, ya que es el único hilo que lee del dispositivo), acumula los bytes recibidos, localiza el marcador de trama y emite la señal \texttt{frameReady(QImage, fps)} por cada fotograma completo.
    \item \textbf{Hilo de detección de rostros} (\texttt{detectionThreadFunc}): coge el último fotograma disponible (compartido mediante \texttt{m\_latestFrame} protegido por \texttt{m\_latestFrameMutex}), lanza el clasificador Haar de OpenCV y emite la señal \texttt{facesDetected(QVector<QRect>, detFps)} con los rectángulos detectados.
\end{itemize}

La sincronización entre hilos se apoya en tres objetos distintos, cada una con un propósito concreto:
\begin{itemize}
    \item \texttt{m\_ftdiMutex}: protege únicamente las llamadas a \texttt{FT\_Write} (comandos salientes desde el hilo principal), que pueden solaparse en el tiempo con la lectura continua del hilo de captura sobre el mismo \textit{handle} FTDI.
    \item \texttt{m\_latestFrameMutex}: protege el traspaso del último fotograma capturado desde el hilo de captura hacia el hilo de detección.
    \item \texttt{m\_facesMutex}: protege la lista de rectángulos detectados, consumida por el hilo principal al pintar el fotograma.
\end{itemize}

Cuando no hay datos disponibles (buffer vacío en captura, o detección desactivada), ambos hilos secundarios recurren a una espera fija breve (\texttt{std::this\_thread::sleep\_for}, entre 1 y 30~ms según el caso) en lugar de una variable de condición; es una solución simple, aunque una futura revisión podría sustituirla por primitivas de sincronización basadas en eventos si se necesitara reducir aún más la latencia.

\section{Interfaz gráfica}
\label{sec:gui_software}

La interfaz (\texttt{mainwindow.ui}) se organiza en dos columnas dentro de un \texttt{QHBoxLayout}, tal y como puede verse en la Figura~\ref{fig:GUI1} (Apéndice~\ref{ch:manual_usuario}):

\begin{itemize}
    \item \textbf{Columna izquierda}: la ventana de vídeo (\texttt{lblVideo}), con el indicador de \ac{FPS} de captura (\texttt{lblFpsRx}) y el estado de conexión (\texttt{lblStatus}, en verde si el FT232H está abierto en modo \textit{Sync FIFO} y en rojo si no).
    \item \textbf{Columna derecha}, organizada en \texttt{QGroupBox} independientes: \textit{Cámara} (activar/desactivar captura y alternar entre sensor real e imagen sintética), \textit{Detección de rostros}, \textit{LEDs} (dos botones individuales), \textit{Configuración (I2C)} y \textit{Log de comandos}.
\end{itemize}

Dos botones adicionales (\textit{Enviar Comando RAW\ldots} y \textit{Enviar Comando I2C\ldots}) se añaden de forma \textbf{dinámica} al grupo de \textit{Configuración} desde el propio constructor de \texttt{MainWindow} (no están definidos en el \texttt{.ui}), y abren un \texttt{QInputDialog} para introducir el comando en texto libre.

\section{Protocolo de recepción de imagen}
\label{sec:protocolo_recepcion}

La función \texttt{findMarker} busca en el flujo de bytes acumulado la secuencia \texttt{0xAA 0x55 0xAA 0x55} que marca el inicio de cada fotograma. Una vez localizado el marcador, se descartan los bytes previos y se extraen exactamente \texttt{kFrameBytes} = $640 \times 480 = 307\,200$ bytes de luminancia. A diferencia de reconstruir la imagen RGB duplicando el valor de luminancia en los tres canales (R=G=B=Y), la implementación actual usa directamente el formato nativo \texttt{QImage::Format\_Grayscale8}, evitando la copia y conversión adicional. Si el acumulador de bytes crece por encima de $4\times$ el tamaño de lectura (\texttt{kBufSize} = 65\,536~B) sin encontrar el marcador, se descarta todo el contenido salvo los 3 últimos bytes (para no perder un marcador que quede entre dos lecturas), limitando así el crecimiento de memoria ante una posible pérdida de sincronismo prolongada.

El indicador de \ac{FPS} de captura (\texttt{lblFpsRx}) se calcula con una ventana deslizante de hasta 30 marcas de tiempo (\texttt{std::deque}) sobre los fotogramas efectivamente recibidos, no sobre los bytes leídos del USB.

\section{Protocolo de comandos y control}
\label{sec:protocolo_comandos}

Todos los comandos comparten el mismo byte de sincronismo (\texttt{kCmdSync~=~0xCC}) y se serializan en tramas de 4 o 6 bytes mediante \texttt{sendCmd4}/\texttt{sendCmdI2C}, coherentes con el protocolo de \texttt{ftdi\_controller} descrito en la Sección~\ref{sec:ftdi_controller}. La Tabla~\ref{tab:comandos_sw} resume los identificadores de comando definidos en el código (\texttt{kCmdLed}, \texttt{kCmdBcd}, \texttt{kCmdI2c}, \texttt{kCmdCap}, \texttt{kCmdSrc}), que coinciden con los CMD~0x01–0x05 del \texttt{cmd\_processor} (Sección~\ref{sec:cmd_processor}).

\begin{table}[H]
\centering
\caption{Comandos PC$\rightarrow$FPGA definidos en \texttt{mainwindow.cpp}.}
\label{tab:comandos_sw}
\begin{tabular}{>{\raggedright\arraybackslash}p{2.2cm}>{\raggedright\arraybackslash}p{2.5cm}>{\raggedright\arraybackslash}p{8.3cm}}
\rowcolor[HTML]{156082}
{\color{white}\textbf{Constante}} & {\color{white}\textbf{Valor}} & {\color{white}\textbf{Origen en la GUI}} \\
\rowcolor[HTML]{C0E6F5}
\texttt{kCmdCap} & 0x04 & Botón \textit{Deshabilitar/Habilitar cámara} (\texttt{onToggleCamera}). \\
\texttt{kCmdSrc} & 0x05 & Botón \textit{Fuente: Real / Modo: Sintético} (\texttt{onToggleImageSource}). \\
\rowcolor[HTML]{C0E6F5}
\texttt{kCmdLed} & 0x01 & Botones \textit{LED 0}/\textit{LED 1} (\texttt{onLed0}/\texttt{onLed1}). \\
\texttt{kCmdI2c} & 0x03 & Diálogo \textit{Enviar Comando I2C\ldots} (\texttt{onSendI2CCommand}), trama de 6 bytes con página, registro y dato. \\
\rowcolor[HTML]{C0E6F5}
\end{tabular}
\end{table}

Los botones \textit{LED~0}/\textit{LED~1} envían en el campo de datos una máscara combinada de 2 bits (\texttt{bit0}=estado de LED~0, \texttt{bit1}=estado de LED~1) mediante \texttt{kCmdLed}. Sin embargo, en el procesador de comandos actualmente se interpreta como un \textbf{pulso de conmutación} (\texttt{s\_led\_toggle\_r}) que alterna el LED de la FPGA, no como una escritura directa de una máscara de valor.

Además de los comandos con botón propio, dos diálogos genéricos permiten enviar tramas arbitrarias sin recompilar la aplicación:
\begin{itemize}
    \item \textbf{Comando RAW} (\texttt{onSendRawCommand}): solicita una lista de bytes en hexadecimal o decimal separados por espacios y los transmite tal cual con \texttt{FT\_Write}, validando que cada valor esté en el rango de 8~bits.
    \item \textbf{Comando I2C} (\texttt{onSendI2CCommand}): solicita página, dirección de registro y dato (p.~ej. \texttt{1 0x37 0x0080} para forzar los 30~fps del registro R55 descrito en la Sección~\ref{sec:top}), valida los rangos (página/registro de 8~bits, dato de 16~bits) y compone la trama de 6~bytes correspondiente. Esta es la forma utilizada en la Figura~\ref{fig:GUI2} del Apéndice~\ref{ch:manual_usuario}.
\end{itemize}

\section{Detección de rostros}
\label{sec:deteccion_rostros}

Al arrancar, \texttt{MainWindow} intenta cargar el clasificador Haar \texttt{haarcascade\_frontalface\_default.xml} desde el directorio del ejecutable (\texttt{QCoreApplication::applicationDirPath()}); si falla, la detección queda deshabilitada y el botón correspondiente lo indica en la barra de estado. Cuando está activa, el hilo de detección ejecuta \texttt{cv::CascadeClassifier::detectMultiScale} sobre el último fotograma en escala de grises (factor de escala 1.2, mínimo 4 vecinos, tamaño mínimo de detección $40\times40$~px) y los rectángulos resultantes se dibujan en rojo sobre el \texttt{QPixmap} mostrado, en el hilo principal, dentro de \texttt{onFrameReady}.

\section{Gestión de errores y limitaciones}
\label{sec:limitaciones_sw}

El principal problema de la arquitectura original —el bloqueo de la interfaz gráfica mientras \texttt{FT\_Read} esperaba datos— queda resuelto en la versión actual al ejecutar la captura en su propio hilo y comunicar los resultados a la \ac{GUI} exclusivamente mediante señales de Qt, sin llamadas bloqueantes en el hilo principal. Las limitaciones que persisten son de menor importancia:


\begin{itemize}
    \item \texttt{CMakeLists.txt} fija \texttt{WIN32\_EXECUTABLE TRUE} en Windows, por lo que la aplicación no abre consola: cualquier traza de depuración enviada a \texttt{stdout}/\texttt{stderr} queda oculta.
    \item Si \texttt{FT\_Open} falla al arrancar, la aplicación se limita a mostrar un mensaje único en la barra de estado y continúa sin conexión; no existe una lógica de reintento ni de reconexión en caliente si el cable USB se desconecta y se vuelve a conectar durante la ejecución.
    \item Las esperas de los hilos de captura y detección cuando no hay datos disponibles usan temporizaciones fijas (\texttt{sleep\_for} de 1--30~ms) en lugar de variables de condición; funcional pero mejorable si se quisiera minimizar la latencia al máximo.
\end{itemize}

\end{document}